\chapter{Helicobacter pylori Test}

\section*{What Is This Test?}
This test detects infection by \textit{Helicobacter pylori}, a spiral bacterium that colonizes the stomach lining. H. pylori infection is a major cause of gastritis, peptic ulcer disease, and is associated with gastric cancer \cite{crowe2019helicobacter}.

\section*{Alternative Names}
\begin{itemize}
  \item H. pylori stool-antigen test
  \item H. pylori breath test
  \item Urea breath test (UBT)
  \item H. pylori antibody test
\end{itemize}

\section*{Principle}
Three primary methods are used:
\begin{itemize}
    \item \textbf{Urea Breath Test (UBT):} The patient swallows 13C- or 14C-labeled urea. H. pylori produces urease, which hydrolyzes urea into labeled carbon dioxide (CO2). Exhaled labeled CO2 is detected in breath samples.
    \item \textbf{Stool Antigen Test:} Uses monoclonal antibodies to detect H. pylori antigens in stool samples.
    \item \textbf{Serology:} Detects IgG antibodies against H. pylori in serum, but cannot differentiate active from past infection.
\end{itemize}

\section*{History}
H. pylori was discovered in 1982 by Australian physicians Barry Marshall and Robin Warren, who demonstrated that it causes peptic ulcers. They were awarded the Nobel Prize in Medicine in 2005. Non-invasive breath and stool tests were developed in the 1990s to replace endoscopy for routine screening.

\section*{Why This Test Is Ordered}
This test is ordered if you have symptoms of gastritis or peptic ulcers, such as gnawing or burning stomach pain, bloating, nausea, vomiting, or unexplained weight loss. It is also used to verify that an H. pylori infection has been successfully eradicated.

\section*{Primary vs Secondary Test}
This is a primary test for diagnosing active H. pylori infection.

\section*{How This Test Is Performed}
- **Breath Test:** You blow into a collection bag, drink a labeled urea solution, wait 10-15 minutes, and blow into a second bag.
- **Stool Test:** You collect a stool sample in a clean container.
- **Blood Test:** Standard venipuncture.

\section*{What Preparation Is Needed}
\textbf{Crucial preparation is required:} You must stop taking antibiotics, bismuth subsalicylate (Pepto-Bismol), and proton pump inhibitors (PPIs) for at least 2 weeks before a breath or stool test, as these medications cause false-negative results.

\section*{Contraindications and Precautions}
\begin{itemize}
  \item There are no contraindications for the stool antigen test. The breath test is safe but should be performed with caution in pregnant women if using 14C-urea (which has a tiny dose of radioactivity; 13C-urea is non-radioactive and preferred).
\end{itemize}
\section*{Risks}
There are no risks associated with the breath or stool tests. Standard venipuncture risks apply to the blood test.

\section*{What Happens After the Test}
If the test is positive, your doctor will prescribe "triple therapy" or "quadruple therapy" consisting of two antibiotics (e.g., clarithromycin, amoxicillin) and a PPI to eradicate the bacterium \cite{chey2017treatment}.

\section*{Understanding Results}

\subsection*{Below Normal}
A negative result indicates that H. pylori antigen or urease activity was not detected, suggesting that you do not have an active H. pylori infection.

\subsection*{Above Normal}
A positive stool antigen or breath test confirms an active H. pylori infection. A positive antibody test confirms exposure to the bacterium, but does not prove active infection.

\section*{Normal Results}
Negative.

\section*{Follow-up Tests}
Following eradication therapy, a repeat breath or stool test must be performed at least 4 weeks after completion of treatment to confirm successful eradication of the infection \cite{malfertheiner2017management}.

\section*{References}
\bibliographystyle{unsrtnat}
\bibliography{references}

\clearpage
\section*{Regulatory Status and Authorization Date}
This chapter describes a test category, not a specific commercial product. FDA, CDSCO, EU IVDR, and MHRA approval or registration dates therefore cannot be assigned without the assay manufacturer, product name, instrument, and intended use. For laboratory-developed tests, an individual agency approval date may not exist. See the Preface for the interpretation of regulatory status.

\clearpage
